\documentclass[a4paper,11pt]{article}
\usepackage[utf8]{inputenc}
\usepackage[T1]{fontenc}
\usepackage{geometry}
\usepackage{enumitem}
\usepackage{titlesec}
\usepackage[table]{xcolor} % Combined xcolor and colortbl for better compatibility
\usepackage{fancyhdr}
\usepackage{tikz}
\usepackage{tabularx}
\usepackage{graphicx}
\usepackage{lastpage} % Added for dynamic total page count
\usepackage{hyperref} % Should be loaded last

% Page Setup
\geometry{top=1in, bottom=1in, left=0.8in, right=0.8in}
\pagestyle{fancy}
\fancyhf{}
\lhead{\small Project 83: Agentic AI Compiler Assistant}
\rhead{\small Week 1 Report}
\cfoot{\thepage\ of \pageref{LastPage}}

% Formatting
\titleformat{\section}{\Large\bfseries\color{blue!70!black}}{\thesection.}{1em}{}[\titlerule]
\titleformat{\subsection}{\large\bfseries\color{blue!50!black}}{\thesubsection}{1em}{}

% Custom commands
\newcommand{\statusgood}[1]{\colorbox{green!20}{\textbf{#1}}}
\newcommand{\statuspending}[1]{\colorbox{yellow!20}{\textbf{#1}}}
\newcommand{\statusissue}[1]{\colorbox{red!20}{\textbf{#1}}}

\begin{document}

\begin{center}
    {\huge \textbf{Week 1 Progress Report}} \\
    \vspace{3mm}
    {\large \textbf{Project 83: Conversational Agentic AI Compiler Assistant}} \\
    \vspace{2mm}
    \textit{Problem Definition \& Scope Research} \\
    \vspace{2mm}
    \textbf{Reporting Period:} Week 1 | \textbf{Status:} \statusgood{COMPLETED}
\end{center}

\vspace{5mm}

\section{Executive Summary}

Week 1 focused on establishing the foundational understanding of the project scope, defining the Mini-Python language subset, and designing the agentic interaction flow. This week was critical in setting clear boundaries for the compiler's capabilities and planning the integration of AI-driven assistance. The primary deliverables included a formal EBNF grammar specification, comprehensive documentation of the Mini-Python subset, and a detailed research analysis of agentic AI patterns suitable for compiler assistance.

\textbf{Key Achievement:} Successfully defined a complete EBNF grammar for the Mini-Python subset covering variables, control flow (if/else, while, for), functions, and basic data types. The scope carefully balances educational complexity with implementation feasibility within the 16-week timeline.

\subsection{Week Overview}

\begin{itemize}[leftmargin=*]
    \item \textbf{Planned Objectives:} Define EBNF grammar, map agentic interaction flow, establish project documentation standards
    \item \textbf{Achieved Objectives:} Complete EBNF specification, research ReAct agent patterns, create initial documentation framework
    \item \textbf{Completion Rate:} 100\%
\end{itemize}

\section{Detailed Activities \& Accomplishments}

\subsection{1.1 Mini-Python Language Definition}

\textbf{Objective:} Define a formal grammar specification for the subset of Python to be supported by the compiler.

\textbf{Activities Performed:}
\begin{itemize}[leftmargin=*]
    \item Conducted extensive analysis of Python's official grammar specification from the Python Language Reference
    \item Identified core language features essential for educational compiler development:
    \begin{itemize}
        \item Variable declarations and assignments (integer, string, boolean types)
        \item Arithmetic and logical operators with proper precedence rules
        \item Control structures: if/elif/else conditionals, while loops, for loops
        \item Function definitions with parameters and return statements
        \item Python-specific features: significant indentation, dynamic typing
    \end{itemize}
    \item Developed EBNF grammar notation covering all supported constructs
    \item Created syntax diagrams using LucidChart for visual representation
\end{itemize}

\textbf{Key Decisions:}
\begin{itemize}[leftmargin=*]
    \item \textbf{Excluded Features:} Classes, decorators, exception handling, list comprehensions, lambda functions - these add significant complexity without proportional educational value for a first compiler implementation
    \item \textbf{Type System:} Chose dynamic typing to match Python's philosophy, but with basic type checking in semantic analysis
    \item \textbf{Indentation:} Decided to fully implement Python's significant whitespace, as this is a defining characteristic and valuable learning experience
\end{itemize}

\textbf{Deliverables Produced:}
\begin{itemize}[leftmargin=*]
    \item \texttt{grammar/mini\_python\_grammar.ebnf} - Formal EBNF specification
    \item Grammar visualization diagrams (LucidChart)
    \item Language specification document with examples and test cases
\end{itemize}

\subsection{1.2 Agentic Interaction Flow Design}

\textbf{Objective:} Research and design the AI agent architecture for conversational compiler assistance.

\textbf{Research Activities:}
\begin{itemize}[leftmargin=*]
    \item Studied the ReAct (Reason + Act) framework for building autonomous agents
    \item Analyzed prompt engineering patterns for technical assistance agents
    \item Investigated existing compiler error message systems and their limitations
    \item Explored LLM capabilities for code understanding and debugging assistance
\end{itemize}

\textbf{Design Outcomes:}
\begin{itemize}[leftmargin=*]
    \item \textbf{Multi-Agent Architecture:} Designed a two-agent system:
    \begin{enumerate}
        \item \textbf{Auditor Agent:} Analyzes code for syntax errors, semantic issues, and security vulnerabilities
        \item \textbf{Assistant Agent:} Provides conversational explanations and suggested fixes to users
    \end{enumerate}
    \item \textbf{Interaction Protocol:} Defined the workflow:
    \begin{enumerate}
        \item User submits Python code
        \item Compiler performs traditional analysis (lexing, parsing, semantic analysis)
        \item If errors detected, compiler state (tokens, AST, symbol table) is serialized
        \item Auditor Agent receives compiler state and analyzes issues
        \item Assistant Agent formulates user-friendly explanations and suggestions
        \item User receives conversational feedback with fix recommendations
    \end{enumerate}
\end{itemize}

\textbf{Key Research Findings:}
\begin{itemize}[leftmargin=*]
    \item ReAct pattern allows agents to "think step-by-step" through compiler errors
    \item Providing structured compiler state (AST, symbol table) dramatically improves AI accuracy vs. raw source code only
    \item Conversational debugging reduces user frustration and improves learning outcomes
\end{itemize}

\newpage

\section{Technical Challenges \& Solutions}

\subsection{3.1 Challenge: Indentation Handling in Grammar}

\textbf{Problem:} Python's significant whitespace is notoriously difficult to represent in traditional EBNF grammar, as most grammars ignore whitespace.

\textbf{Solution Approach:}
\begin{itemize}[leftmargin=*]
    \item Researched Python's actual lexer implementation in CPython
    \item Decided to handle indentation at the lexical analysis stage (Lexer) rather than parser
    \item Designed a special token-generation system that emits \texttt{INDENT} and \texttt{DEDENT} tokens
    \item Documented this approach with explicit examples in the grammar specification
\end{itemize}

\textbf{Impact:} This decision provides a clear implementation path for Week 5 (Lexer development) and aligns with industry-standard compiler design patterns.

\subsection{3.2 Challenge: Scope Limitation Balance}

\textbf{Problem:} Determining which Python features to include vs. exclude to maintain educational value while ensuring project feasibility.

\textbf{Solution Approach:}
\begin{itemize}[leftmargin=*]
    \item Created a feature prioritization matrix based on:
    \begin{itemize}
        \item Educational value (fundamental compiler concepts demonstrated)
        \item Implementation complexity
        \item Time constraints (16-week timeline)
    \end{itemize}
    \item Validated scope with compiler design literature
    \item Ensured included features cover all major compiler phases (lexing, parsing, semantic analysis, basic optimization opportunities)
\end{itemize}

\textbf{Impact:} Final scope provides sufficient complexity for comprehensive learning while remaining achievable within timeline constraints.

\section{Documentation \& Knowledge Management}

\subsection{4.1 Documentation Framework Established}

Created comprehensive documentation structure:
\begin{itemize}[leftmargin=*]
    \item \texttt{README.md} - Project overview and quick start guide
    \item \texttt{docs/DetailedExecutionPlanandReport.tex} - 16-week execution roadmap
    \item \texttt{grammar/} - Formal language specifications
    \item \texttt{docs/prompts.md} - AI prompt templates and system instructions
\end{itemize}

\subsection{4.2 Tools \& Resources Utilized}

\begin{table}[h]
\centering
\begin{tabularx}{\textwidth}{|l|X|}
\hline
\rowcolor{blue!10}
\textbf{Tool/Resource} & \textbf{Purpose} \\
\hline
LucidChart & Grammar visualization and architecture diagrams \\
\hline
Python Language Reference & Authoritative grammar specification \\
\hline
Markdown & Documentation and project planning \\
\hline
ReAct Framework Papers & Agentic AI design patterns \\
\hline
Dragon Book (Compiler Textbook) & Theoretical foundations \\
\hline
\end{tabularx}
\caption{Tools and Resources - Week 1}
\end{table}

\section{Metrics \& Progress Indicators}

\begin{table}[h]
\centering
\begin{tabularx}{\textwidth}{|l|c|X|}
\hline
\rowcolor{blue!10}
\textbf{Metric} & \textbf{Target} & \textbf{Actual} \\
\hline
EBNF Grammar Completeness & 100\% & 100\% - All core constructs defined \\
\hline
Documentation Quality & High & High - Comprehensive with examples \\
\hline
Research Depth & Thorough & Thorough - Multiple sources consulted \\
\hline
Scope Clarity & Clear & Clear - Well-defined boundaries \\
\hline
Schedule Adherence & On Track & On Track - Week 1 completed as planned \\
\hline
\end{tabularx}
\caption{Week 1 Performance Metrics}
\end{table}

\section{Risk Assessment \& Mitigation}

\subsection{6.1 Identified Risks}

\begin{enumerate}
    \item \textbf{Grammar Complexity Risk (Low):} Initial grammar may require refinement during implementation
    \begin{itemize}
        \item \textit{Mitigation:} Documented grammar with extensive examples; prepared to iterate in Week 5-6
    \end{itemize}
    
    \item \textbf{AI Capability Risk (Medium):} Uncertainty about Gemini API's ability to handle complex AST structures
    \begin{itemize}
        \item \textit{Mitigation:} Week 2 dedicated to hands-on API testing to validate capabilities early
    \end{itemize}
    
    \item \textbf{Scope Creep Risk (Medium):} Temptation to add features beyond defined scope
    \begin{itemize}
        \item \textit{Mitigation:} Formal scope document with clear exclusions; strict adherence to EBNF grammar
    \end{itemize}
\end{enumerate}

\section{Lessons Learned}

\begin{itemize}[leftmargin=*]
    \item \textbf{Early Planning Value:} Investing time in comprehensive grammar definition prevents ambiguity in later implementation phases
    \item \textbf{Research Importance:} Understanding ReAct patterns early provides clear direction for AI integration design
    \item \textbf{Documentation Discipline:} Establishing documentation standards from Week 1 creates good habits for the entire project
    \item \textbf{Realistic Scoping:} Being conservative with feature scope reduces risk of timeline overruns
\end{itemize}

\section{Next Week Preview (Week 2)}

\textbf{Focus:} Gemini API \& Agentic Workflow Design

\textbf{Planned Activities:}
\begin{itemize}[leftmargin=*]
    \item Set up Google AI Studio development environment
    \item Obtain and configure Gemini API credentials
    \item Conduct capability testing with raw Python tokens and simple AST structures
    \item Design detailed prompt architecture for Auditor and Assistant agents
    \item Create prompt templates with EBNF grammar injection
    \item Test prompt effectiveness with sample error scenarios
\end{itemize}

\textbf{Success Criteria:}
\begin{itemize}[leftmargin=*]
    \item Successful API connection and authentication
    \item Demonstrated ability of Gemini to interpret compiler state information
    \item Functional prompt templates that produce helpful, accurate responses
\end{itemize}

\section{Conclusion}

Week 1 successfully established the foundational understanding necessary for the Agentic AI Compiler Assistant project. The formal EBNF grammar provides a clear specification for implementation, while the agentic architecture design offers an innovative approach to compiler error reporting. All planned objectives were achieved on schedule with high-quality deliverables.

The project is well-positioned to proceed to Week 2's hands-on API exploration and prompt engineering work. The clear scope definition and thorough research conducted this week will pay dividends throughout the remaining 15 weeks of development.

\vspace{5mm}
\noindent\textbf{Prepared by:} Project Team \\
\textbf{Reporting Date:} Week 1 Completion \\
\textbf{Next Review:} Week 2 Report

\end{document}